\chapter{Choice Before, During, and After}\label{ch:choice}

Selective cinema is often described as giving the viewer a choice, and the description hides several different arrangements. A viewer may choose once, before the film, or repeatedly during it. The choice may be the viewer's own, delegated to a companion, left to chance, or made by a system on the viewer's behalf. Each arrangement produces a different work and a different relation between the work and the person watching it. This chapter distinguishes them and argues against the assumption that more choice is always better.

\section{Preselection}

The simplest arrangement is a choice made before the film begins and kept throughout. The viewer sets the dial, perhaps on the advice of a description of each realization, and watches one realization from start to finish.

Preselection asks nothing of the switch system, because nobody switches. It gives each viewer a complete, coherent film. It makes the choice deliberate, a statement of preference made with some knowledge of the options. Its limitation is that the choice is made before the viewer knows what the film is like. A viewer who chose the thriller and finds it too much, or the beginner realization and finds it too slow, is held to a decision made in ignorance.

\section{Switching during the film}

Allowing viewers to switch during the film turns the family into something that can be explored. A viewer can sample realizations, move toward more explanation during a hard stretch, or follow a character who has become more interesting. The family stops being a set of films and becomes a space the viewer moves through.

Switching is also where the switch system of \Cref{ch:switch-admissibility} does its work. A viewer who switches carries a history, and not every switch keeps that history coherent. The exhibition has three broad options. It can \emph{lock} the dial except for admissible switches, so that the viewer can only make moves the family supports. It can \emph{warn}, allowing any switch but signaling when a move will leave the viewer without something they need or contradict something they were shown. Or it can restrict switching to \emph{rendezvous}, where, for viewers who have not yet switched, every move is admissible and nothing needs to be tracked.

Each option expresses a different view of the viewer. Locking treats the family's coherence as the artist's responsibility and protects the viewer from incoherence they did not intend. Warning treats the viewer as entitled to make incoherent moves knowingly, perhaps out of curiosity about what the other realization is doing. Rendezvous-only switching is the simplest, needs no state in the glasses, and gives the film a rhythm of choice points like chapter breaks.

\section{The artist's right to constrain}

It is tempting to regard any restriction on switching as a limitation to be minimized, and to treat the ideal family as one in which every viewer can move anywhere at any time. The book does not take that view.

A novel does not let its reader substitute paragraphs. A symphony does not let its listeners change the order of movements. These are not defects. They are part of what makes the work a work: a designed sequence, whose effects depend on its order. Music offers a closer precedent than either. Eco described ``open works'' whose performers or readers complete them within limits the author sets, such as Stockhausen's \emph{Klavierstück XI}, whose pianist chooses the order of its fragments during the performance \cite{eco}. The freedom in such works is real and it is composed. A variant family is a designed space of the same kind rather than a designed sequence, but it is still designed, and its effects can depend on where and when a viewer is allowed to move. A mystery whose realizations follow different investigators may depend on each viewer staying with one investigator until the solution. A family of endings may depend on the choice of ending being made at the branch and not before.

An artist who restricts switching is therefore not failing to offer choice. They are deciding what kind of work to make. The switch system makes such restrictions principled rather than arbitrary: they can be tied to coherence, allowing moves that keep the viewer's history intact and forbidding those that would break it, rather than simply closing the dial. But even a restriction that goes beyond coherence, keeping a viewer on one realization for reasons of form, is a legitimate artistic choice, provided it is declared. A viewer should know, before choosing, whether they will be able to change their mind.

\section{Delegation}

A viewer may let someone else choose. A parent may choose for a child, a friend may choose for a friend, a teacher for a class. Delegation is common and often reasonable. A young child may not be able to read a description of the realizations, and a parent's choice is the ordinary way such decisions are made.

It is worth noticing how much the complementary design of \Cref{ch:genre} eases delegation. When the realizations are the thriller and the children's satire of the \emph{Protocols} pair, a parent choosing the satire for a child is choosing an address, not guarding against harm, and the stakes of a wrong choice are low. When the realizations include one that a child should not see, delegation carries a burden the glasses cannot share, because, as \Cref{ch:unfiltered} showed, the child can remove the glasses and see the composite. Delegation in a restrictive family is a matter of trust and supervision, not of technology.

\section{Random assignment}

A screening can assign realizations at random: each ticket, each pair of glasses, each seat, set to a realization nobody chose. This makes variation part of the event. The audience knows the family exists and does not know which realization each of them will receive, and the conversation afterward begins from genuine discovery.

Random assignment is a legitimate artistic device, with one condition carried over from \Cref{ch:after-screening}: the existence of variation must be disclosed. Viewers can agree to be surprised by which realization they receive. They cannot meaningfully agree to a variation they do not know is happening.

\section{Adaptive selection}

The last arrangement is the one this book most resists. A system could choose the realization for the viewer, and change it during the film, on the basis of what it observes: where the viewer looks, how they react, how their heart rate changes, what they have watched before, what is known about them. Such a system would promise a film continuously tailored to its viewer.

Adaptive selection requires the glasses, or something in the room, to observe the viewer. It would extend into the cinema the practice, well documented in digital services, of collecting behavioral data in order to predict and shape what people do \cite{zuboff}. It abandons the architectural guarantee of \Cref{ch:phase}, that the simplest glasses cannot record or measure anything, and replaces it with a sensor pointed at a person. It moves the choice from the viewer to a system whose reasoning the viewer cannot see. And it makes the question \Cref{ch:depth} raised about expertise, who decides what a viewer can follow, into a question answered continuously by inference, without the viewer's participation.

None of this is inherent in selective cinema. The glasses of Part~I are timing devices. They can be set by hand, and they need to know nothing about the person wearing them. Adaptive selection is an addition, and it should be treated as one: never the default, always declared, and always refusable in favor of a realization the viewer chooses. \Cref{ch:persuasion} returns to the specific danger of a system that observes viewers and then chooses what they see.

\section{Choosing to choose}

The arrangements in this chapter are not exclusive, and a single screening can combine them: a preselected realization with switching allowed at rendezvous, or random assignment with a warning dial for the curious. What matters is that the arrangement is part of the work and is known to the audience. A viewer entering a selective screening should know three things before it begins: what the realizations are, how their own will be chosen, and whether they will be able to change it. With those three things known, choice before, during, and after the film becomes one more element an artist can compose, and one the audience can understand.
